% !PRINTED 278
\weberpage{278}
%<<CONT
ist, so setzen wir
\weq{14}{U = uvv'v'', \quad U_1 = u_1v_1v_1'v_1'', \quad U_2 = u_2v_2v_2'v_2'',}
\weq{15}{
\begin{array}{rl}
2A & = U + (-1)^\mu(U_1+U_2), \\
B & = UU_1 + UU_2 + (-1)^\mu U_1U_2, \\
& = \dfrac{4}{U_1} + \dfrac{4}{U_2} + (-1)^\mu\dfrac{4}{U}.
\end{array}
}

Aus (12) und § 73, (4) erhalten wir dann folgende zusammengehörige
Vertauschungen:
\weq{16}{
\begin{array}{lll}
\omega, & A, & B, \\
-\dfrac{1}{\omega}, & \varrho^{n'+1}A, & \varrho^{-(n'+1)}B, \\
\omega+1, & \sigma^{n'+1}e^{\frac{2\mu\pi i}{3}}A, & \sigma^{-(n'+1)}e^{-\frac{2\mu\pi i}{3}}B.
\end{array}
}

Wir wenden unser Prinzip zur Herleitung von Modulargleichungen auf
diese Funktionen an und bemerken dazu noch folgendes:

Die in (16) vorkommenden dritten Einheitswurzeln sind $=1$, wenn
entweder $n$ durch 3 unteilbar und $\mu$ durch 3 teilbar ist, oder $n''$
durch 3 teilbar ist und $n'$ den Rest 2 läßt. In diesen Fällen ist jede
rationale Funktion von $A$ und $B$ Wurzel einer Transformationsgleichung.
In den anderen Fällen kommt diese Eigenschaft dem Kubus einer solchen
rationalen Funktion von $A$ und $B$ zu, bei denen die Differenzen der
Exponenten sämtlicher Glieder einander nach dem Modul 3 kongruent sind.

Sind diese rationalen Funktionen \gesperrt{ganze} Funktionen und sind
außerdem ihre sämtlichen Werte für $q = 0$ endlich, so müssen sie einer
Konstanten gleich sein, und dadurch gewinnen wir Transformationsgleichungen.

Sind $n', n''$ Primzahlen, so genügt es auch hier (vgl. § 72), wenn die
negativen Potenzen von $q$ in der Entwickelung einer solchen Funktion
nach steigenden Potenzen von $q$ in dem einen Hauptfall wegfallen,
nämlich in dem, wo
\[
\begin{pmatrix} a, & 0 \\ c, & \partial \end{pmatrix}, \quad
\begin{pmatrix} a', & 0 \\ c', & \partial' \end{pmatrix}, \quad
\begin{pmatrix} a'', & 0 \\ c'', & \partial'' \end{pmatrix}
\]
gleich sind
\[
\begin{pmatrix} 1, & 0 \\ 0, & n \end{pmatrix}, \quad
\begin{pmatrix} 1, & 0 \\ 0, & n' \end{pmatrix}, \quad
\begin{pmatrix} 1, & 0 \\ 0, & n'' \end{pmatrix},
\]
also
\[
U = f(\omega)f(n'\omega)f(n''\omega)f(n'n''\omega);
\]
%CONT>>
